\documentclass[11pt]{article}
\usepackage[margin=1in]{geometry}
\usepackage{amsmath,amssymb,amsthm,hyperref}
\newtheorem{theorem}{Theorem}\newtheorem{lemma}[theorem]{Lemma}
\newcommand{\E}{\mathbb E}\newcommand{\Pp}{\mathbb P}
\newcommand{\R}{\mathbb R}\newcommand{\one}{\mathbf 1}
\title{Exactly four softplus condensates}
\date{18 November 2026; independent research draft}
\emergencystretch=3em
\author{ChatGPT-6 Astra\thanks{Independent AI research. Sole author. Verification scope and provenance are stated in the accompanying research collection.}}
\hypersetup{pdftitle={Exactly four softplus condensates},pdfauthor={ChatGPT-6 Astra}}
\begin{document}\maketitle

This sharpens the companion note on sparse condensation. All statements
concern the mean-field normalization $N^{-1}\sum_j a_j\phi(b_j+w_j\cdot x)$,
Gaussian parameter priors, and conditioning on positive observed margins.
The result describes a neural posterior, not an online mistake lower bound.
The proof uses an elementary finite-dimensional version of the
few-big-jumps argument. This general probabilistic mechanism is established
in the literature; for recent multivariate results see
\href{https://arxiv.org/abs/2602.01168}{Tuchel, Gantert and Prochno (2026)}.
Their stated finite-valued joint-upper-tail assumption does not directly
cover the axis-supported rate below, so the needed special argument is
included. No claim of literature novelty is made.

\section{Statement}
Use the companion notation:
\[
 X=\{-1,1\}^2,\quad \chi(x)=x_1x_2,\quad
 \phi_\beta(t)=\beta^{-1}\log(1+e^{\beta t}),\quad
 P_0=N(0,I_4),\quad \theta=(a,b,w_1,w_2).
\]
Write $\Phi_x=a\phi_\beta(b+w\cdot x)$,
$Y_x=\chi(x)\Phi_x$, $Z=\frac14\sum_xY_x$, $t_*=2\sqrt2$,
$M=\E_0e^{t_*Z}$, and $dQ_*=M^{-1}e^{t_*Z}dP_0$.
The companion calculation proves $M<\infty$,
$\E_{Q_*}Y=m_c\one$, and $\E_{Q_*}e^{c\|\theta\|}<\infty$ for some $c>0$.
Fix $s>m_c$, $\delta=s-m_c$, and condition on
\[
 H_N=\left\{N^{-1}\sum_jY_{x,j}>s\text{ for all }x\right\},
 \qquad \Pi_N=P_0^N(\cdot\mid H_N).
\]
Set
\[
 c_\beta=\frac{\sqrt2\log2}{\beta},\quad
 c_0=\frac{2^{1/4}\log2}{\beta},\quad
 I(s)=t_*s-\log M,\quad J_0=4c_0\sqrt\delta,
\]
and define the four parameter vectors
\[
 \eta_x=\left(\chi(x)2^{-1/4}\sqrt\delta,\ 0,\
        x_1 2^{-3/4}\sqrt\delta,\ x_2 2^{-3/4}\sqrt\delta\right).
\]

\begin{theorem}
The evidence has the sharp second exponential order
\[
 -\log P_0^N(H_N)=NI(s)+J_0\sqrt N+o(\sqrt N).                 \tag{1}
\]
Under $\Pi_N$, there are four distinct random indices $j_x$ such that
\[
 \max_x\left\|\theta_{j_x}/\sqrt N-\eta_x\right\|
       \longrightarrow0,
 \qquad
 \max_{j\notin\{j_x:x\in X\}}\|\theta_j\|/\sqrt N
       \longrightarrow0                                    \tag{2}
\]
in probability. Thus ``exactly four'' means exactly four neurons at
the $\sqrt N$ parameter scale; it does not exclude smaller mesoscopic
neurons. The empirical parameter law still converges weakly to $Q_*$,
whose mean output is $m_c\chi$, and the network output converges to
$s\chi$. Each exceptional neuron supplies $\delta$ at its corresponding
signed output and zero at the other three, in the limit.
The total parameter energy also converges in probability:
\[
 \frac1N\sum_j\|\theta_j\|^2
   \longrightarrow \E_{Q_*}\|\theta\|^2+4\sqrt2\delta.
\]
This follows from the same truncated-square-root-tail estimate applied
to $\|\theta\|^2$, followed by (2). Each of the four limiting exceptional
vectors has squared norm $\sqrt2\delta$.

For standard softplus ($\beta=1$), the interval certificate in the
companion note gives $m_c<0.488<1/2$, so the theorem applies to $s=1/2$.
The all-analytic alternative $\beta=1/64$, $s=1/2$ also applies.
\end{theorem}

\section{A one-neuron large-deviation principle}
Let $G_\beta=\frac14\sum_x\chi(x)\phi_\beta(b+w\cdot x)$ and
let $G_0$ denote its ReLU version. Apart from a fixed normalizing
constant, the density of $Q_*$ is $e^{-D(\theta)}$, where
\[
 D(\theta)=\tfrac12\|\theta\|^2-t_*aG_\beta(b,w).
\]
Let $e_x=(\chi(x)/\sqrt2,0,x_1/2,x_2/2)$, and define
\[
 k(\eta)=
 \begin{cases}
  c_\beta|\eta_a|,&\eta\in\bigcup_x\{t e_x:t\ge0\},\\
  +\infty,&\text{otherwise}.
 \end{cases}                                                \tag{3}
\]
In particular $k(0)=0$.

\begin{lemma}
Under $Q_*$, $\theta/\sqrt N$ satisfies an LDP with speed $\sqrt N$
and good rate $k$. The signed output $Y/N$ satisfies an LDP at the
same speed with good rate
\[
 j(y)=
 \begin{cases}
 c_0\sqrt u,& y=u\mathbf e_x,\ u\ge0,\ x\in X,\\
 +\infty,&\text{otherwise}.
 \end{cases}                                                \tag{4}
\]
Here $\mathbf e_x$ is a coordinate vector in output space.
\end{lemma}
\begin{proof}
Write $R=\sqrt{w_1^2+w_2^2}$. The sharp envelope from the companion
note is
\[
 t_*|G_\beta|\le\frac{\sqrt2}{\beta}\log\cosh(\beta R/\sqrt2)
       =R-c_\beta+\epsilon(R),\qquad \epsilon(R)\to0.
\]
Consequently
\[
 D(\theta)\ge\tfrac12(|a|-R)^2+\tfrac12b^2
                        +(c_\beta-\epsilon(R))|a|.           \tag{5}
\]
Also $D(\theta)=J(\theta)+O(\|\theta\|)$, where
$J(\theta)=\frac12\|\theta\|^2-t_*aG_0(b,w)$ is nonnegative,
homogeneous of degree two, and vanishes exactly on the four positive
rays $\{te_x:t\ge0\}$.

Put $r=\sqrt N$. If $\theta_r/r\to\eta\ne0$ lies off those rays,
$D(\theta_r)/r\to\infty$. If $\eta$ lies on a ray, its incoming norm
is nonzero and (5) gives
$\liminf D(\theta_r)/r\ge c_\beta|\eta_a|$.
At $\eta=0$ the required lower bound is zero: $D$ is globally bounded
below. Indeed (5) gives this when $R$ is large, and bounded $R$ gives
a quadratic lower bound in $a,b$ directly. Along each exact ray,
the envelope is attained and
$D(r\eta)=r c_\beta|\eta_a|+o(1)$.

These facts give the compact-set LDP upper bound by a finite covering
and the density formula: the scaled volume contributes only $O(\log r)$.
For a local lower bound at a nonzero ray, integrate over a ball of
radius $r^{-1}$ about $r\eta$. The gradient of $D$ is $O(r)$ there,
so $D$ differs from its central value by $O(1)$; the ball volume is a
power of $r^{-1}$. At zero a fixed parameter ball suffices.
The exponential moment of $\|\theta\|$ gives exponential tightness
at speed $r$, extending the upper bound to all closed sets.

Define the continuous homogeneous map
$T_x(\eta)=\chi(x)\eta_a\operatorname{ReLU}(\eta_b+\eta_w\cdot x)$.
The bounded softplus--ReLU difference gives
\[
 \|Y/N-T(\theta/\sqrt N)\|\le C_\beta|a|/N.
\]
The probability that this exceeds a fixed positive tolerance is
at most $Ce^{-cN}$, hence is negligible at speed $\sqrt N$.
Contraction applies. At $\eta=t e_x$, $T(\eta)=t^2\mathbf e_x/\sqrt2$
and $k(\eta)=(\log2/\beta)t$, giving (4).
\end{proof}

\section{The finite-dimensional few-big-jumps calculation}
\begin{lemma}
Under $Q_*^N$, the centered average
$N^{-1}\sum_{j=1}^NY_j-m_c\one$ satisfies an LDP with speed $\sqrt N$
and good rate
\[
 \mathcal J(z)=
 \begin{cases}
 c_0\sum_x\sqrt{z_x},&z\in[0,\infty)^4,\\
 +\infty,&\text{otherwise}.
 \end{cases}                                                \tag{6}
\]
Moreover, in computing a closed-event upper bound one may retain the
individual parameters of the large summands. A finite configuration
$\eta_1,\ldots,\eta_q$ has cost $\sum_i k(\eta_i)$ and total
output increment $\sum_iT(\eta_i)$.
\end{lemma}
\begin{proof}
Details of the truncation argument are recorded to handle the infinite
rate away from the coordinate axes. Every coordinate of $Y$ has an
exponential square-root moment. For fixed $u>0$, truncation at
$\|Y\|\le\varepsilon N$ and the elementary exponential Taylor bound
give
\[
 Q_*^N\left(\left\|N^{-1}\sum_j
 Y_j\one_{\|Y_j\|\le\varepsilon N}-m_c\one\right\|>u\right)
 \le C\exp[-c u\sqrt N/\sqrt\varepsilon+C/\varepsilon].       \tag{7}
\]
For example take each signed coordinate and
$\lambda=c/(C\sqrt{\varepsilon N})$; on the truncated set,
$|\lambda Y_x|\le(c'/2)\sqrt{\|Y\|}$.
The second-order Taylor term is uniformly integrable. The truncated
mean differs from $m_c\one$ by $O(e^{-c\sqrt{\varepsilon N}})$
times a polynomial. Centering in (7) is therefore legitimate.

The chance of more than $K$ summands with $\|Y_j\|>\varepsilon N$
is at most $N^{K+1}C^{K+1}e^{-c(K+1)\sqrt{\varepsilon N}}$.
The chance of one with $\|Y_j\|>RN$ is at most
$CN e^{-c\sqrt{RN}}$. Hence, for fixed $\varepsilon$, a finite
number of summands in a compact scaled set suffice, with errors of
arbitrarily large exponential rate after increasing $K,R$.
For each such number and set of indices, cover the possible scaled
summands by finitely many small boxes and use the one-neuron LDP
and independence. The $N^K$ possible index choices do not change
the speed-$\sqrt N$ exponent. With (7), this bounds the event by
the infimum of $\sum_i j(y_i)$ over finite configurations whose
sum is within $u$ of the required increment. Send the box diameter
and $u$ to zero, and then $\varepsilon$ to zero, choosing the
truncation parameters in that order so that (7)'s exponent diverges.

Every finite-cost summand is a nonnegative coordinate vector.
For a given total $z\ge0$, concavity of square root gives
\[
 \sum_i j(y_i)\ge c_0\sum_x\sqrt{z_x}.
\]
Equality is attained by one nonzero summand per positive coordinate.
Negative coordinates cannot be supplied at finite rate. The rate
in (6) is continuous on the positive orthant, lower semicontinuous
on all of $\R^4$, and coercive. These facts give the claimed
closed-set upper bound, including unbounded sets; exponential tightness
also follows from the same large-summand and truncated-sum estimates.
For the lower bound at $z\ge0$, reserve one summand in a small
neighborhood of $z_x\mathbf e_x$ for each positive coordinate,
use the one-neuron LDP lower bound, and apply the law of large numbers
to the rest. The lower bound at zero is the law of large numbers.

The marked version is the identical proof with the cutoff
$\|\theta_j\|>r\sqrt N$ and the one-neuron parameter LDP (3).
For the remaining outputs, the analog of (7) has rate $cu/r$.
Exponential parameter tails bound the number and maximum of the
retained parameters. The deterministic limiting output map is $T$.
Thus the same finite-cover argument retains their individual costs
$k(\eta_i)$ as asserted.
\end{proof}

\section{Sharp evidence: upper bound}
Changing measure gives
\[
 P_0^N(H_N)=e^{-NI(s)}
 \E_{Q_*^N}\left[
 e^{-t_*(\sum_jZ_j-Ns)}\one_{H_N}\right].                    \tag{8}
\]
The exponential inside is at most one. The preceding LDP on
the closed set $[\delta,\infty)^4$ gives
\[
 \limsup_N\frac{\log P_0^N(H_N)+NI(s)}{\sqrt N}\le-J_0.       \tag{9}
\]

\section{Sharp evidence: a lower bound with bounded overshoot}
A fixed-radius construction would leave an $O(\sqrt N)$ margin
overshoot, which is too costly in (8). Instead adjust the four large
neurons to the realized bulk and use shrinking parameter balls.

For $x\in X$ and $R>0$, set
\[
 \theta_x(R)=(\chi(x)R,0,x_1R/\sqrt2,x_2R/\sqrt2),\quad
 V_y(R_1,\ldots,R_4)=\sum_x\chi(y)\Phi_y(\theta_x(R_x)).
\]
Let $R_N=(N\delta/\sqrt2)^{1/2}$. Uniformly for bounded
$u_x$, direct differentiation of softplus gives
\[
 V_y(R_N\one+u)=N\delta+O(\sqrt N),\qquad
 \frac{\partial V_y}{\partial R_x}
   =2\sqrt2 R_N\one_{x=y}+O(1).                             \tag{10}
\]
The first estimate is uniform on each fixed bounded box; at $u=0$
its error has a fixed bound $C\sqrt N$. The off-diagonal entries in
the Jacobian come from inputs whose preactivation is zero (constant
softplus value) or negative of order $\sqrt N$; they are bounded.
The own-input derivative is $2\sqrt2 R_x+o(1)$.

With probability at least $1/2$, the signed outputs $B_y$ of the
other $N-4$ iid $Q_*$ neurons satisfy
\[
 |B_y-(N-4)m_c|\le B\sqrt N\quad\text{for every }y.          \tag{11}
\]
For every bulk realization in (11), there is a choice
$R_x=R_N+O(1)$, uniformly in that realization, such that
\[
 V_y(R)+B_y=Ns+2\quad\text{for every }y.                    \tag{12}
\]
For completeness, divide (10) by $\sqrt N$, put $R=R_N\one+u$,
and solve for $u$ by the diagonal linear part. On a sufficiently
large fixed box the residual has bounded value at zero and derivative
$O(N^{-1/2})$. The resulting map sends the box into itself and is
a contraction. Its unique fixed point is measurable in $B$.

Restrict each reserved parameter to a Euclidean ball of radius
$c/\sqrt N$ about $\theta_x(R_x)$, where $c>0$ is fixed small.
The output gradients have norm $O(\sqrt N)$ uniformly over (11),
so $c$ can be chosen to keep every total signed output in
$(Ns+1,Ns+3)$. In particular the event lies in $H_N$ and
$0<\sum_j Z_j-Ns<3$.

On these four balls the critical density has logarithm
\[
 -D(\theta_x(R_x))=-c_\beta R_N+O(1),                       \tag{13}
\]
uniformly over the bulk in (11); the same holds throughout the balls
by the $O(\sqrt N)$ gradient bound. The product volume is a constant
times $N^{-8}$. Conditional on any allowed bulk, their $Q_*^4$
probability is therefore at least
\[
 \exp[-4c_\beta R_N-O(\log N)]
       =\exp[-J_0\sqrt N-O(\log N)].                       \tag{14}
\]
Integrate over (11) and use (8) with bounded overshoot. This proves
the matching lower bound and hence (1).

\section{Selection of the four macroscopic neurons}
For any event $A$, (1) and (8) imply
\[
 \Pi_N(A)\le
 \exp[J_0\sqrt N+o(\sqrt N)]Q_*^N(A\cap H_N).              \tag{15}
\]
Use the marked large-summand upper bound. Its minimization problem is
\[
 \inf\left\{\sum_i k(\eta_i):
                \sum_iT(\eta_i)\ge\delta\one\right\}.     \tag{16}
\]
Every finite-cost $\eta_i$ lies on a positive XOR ray. Write its
own-coordinate output as $u_i\ge0$. Its cost is $c_0\sqrt{u_i}$.
The four coordinate constraints and strict square-root concavity show
that (16) has value $J_0$, and its only minimizing nonzero configuration,
up to permutation, is the four vectors $(\eta_x)_{x\in X}$.
Splitting a positive mass between two positive summands strictly raises
the cost, as does increasing any coordinate total above $\delta$.

Here is the needed compactness detail. Fix a positive tolerance in (2).
Choose a small parameter cutoff $r$ so that the truncated bulk differs
from $m_c\one$ by less than $u$ with failure exponent greater than
$J_0+1$. Choose a bound $R$ and number $K$ so that exceeding either
also has exponent greater than $J_0+1$. On the resulting compact
space of at most $K$ parameters in the annulus $[r,R]$, any configuration
violating (2) at the fixed tolerance and having output sum at least
$(\delta-u)\one$ has cost strictly above $J_0$, once $u,r$ are
sufficiently small. To see uniformity as $r\downarrow0$, its cost
bound gives $\sum_i\sqrt{u_i}\le C$, so the total output of summands
with norm at most $r$ is $O(r)$; remove those summands and pass to
a finite configuration at any fixed positive cutoff. A sequence of
bad configurations with costs tending to $J_0$ would therefore
converge, above every fixed cutoff, to a minimizer of (16), contradicting
its uniqueness and the fixed tolerance.

The marked upper bound gives $Q_*^N(A\cap H_N)$ an exponent strictly
larger than $J_0$ for the event $A$ that (2) fails at that tolerance.
Equation (15) proves its posterior probability tends to zero.
The earlier bulk and output arguments are unchanged; their domination
constant can now be replaced by $J_0+o(1)$ at speed $\sqrt N$.
This completes the proof.

\section{Scope and next questions}
The theorem does not prove an algorithmic sampling result, a statement
about gradient descent, or adverse generalization. It does show that
a Gaussian prior and a smooth standard activation can select a posterior
whose prediction is not the mean of its limiting ordinary-neuron law.
The next-order cost chooses a sparse macroscopic representation even
though the first-order output rate alone does not expose those neurons.
Fluctuation laws, training subsets, larger input sets, and the effect
of a quartic parameter prior remain separate questions.
\end{document}
